% Generated mechanically from frozen input p12/ko/guide.tex.
% Do not edit this generated file; edit the bound locale source instead.

% OK, start here.
%

\setcounter{chapter}{111}
\chapter{문헌 안내}


\phantomsection
\label{guide-section-phantom}


\section{짧은 입문 논문}
\label{guide-section-short-introductions}


\begin{itemize}
\item Barbara Fantechi: \emph{모두를 위한 스택} \cite{fantechi_stacks}
\item Dan Edidin: \emph{스택이란 무엇인가?} \cite{edidin_whatis}
\item Dan Edidin: \emph{곡선의 모듈라이 공간 구성에 관한 강의록}
\cite{edidin_notes}
\item Angelo Vistoli: \emph{대수적 스택과 그 모듈라이 공간의 교차 이론},
특히 부록. \cite{vistoli_intersection}
\end{itemize}


\section{고전적 참고문헌}
\label{guide-section-classics}

\begin{itemize}
\item Mumford: \emph{모듈라이 문제의 피카르 군}
\cite{mumford_picard}
\begin{quote}
Mumford는 이 논문에서 “스택”이라는 용어를 전혀 사용하지 않지만 그 개념은
암묵적으로 들어 있다. 그는 타원곡선의 모듈라이 스택의 피카르 군을 계산한다.
\end{quote}
\item Deligne, Mumford: \emph{주어진 종수를 갖는 곡선 공간의 기약성}
\cite{DM}
\begin{quote}
이 영향력 있는 논문은 오늘날 보편적으로 Deligne--Mumford 스택이라 불리는
의미의 “대수적 스택”을 도입한다. 이는 표현가능한 대각사상을 가지며
스킴에 의한 에탈 표현을 허용하는 스택이다. 많은 기초 결과가
\emph{증명 없이} 제시된다. 이 논문은 스택을 이용하여 종수 $g$인 곡선의
모듈라이 공간이 기약이라는 사실을 두 가지 방식으로 증명한다.
\end{quote}
\item Artin: \emph{버설 변형과 대수적 스택}
\cite{ArtinVersal}
\begin{quote}
이 논문은 Deligne--Mumford 스택을 일반화하는 “대수적 스택”을 도입한다.
오늘날에는 흔히 이를 \emph{Artin 스택}, 즉 표현가능한 대각사상을 가지며
스킴에 의한 매끄러운 표현을 허용하는 스택이라 부른다. 또한 주어진 모듈라이
스택의 표현을 명시적으로 제시하지 않고도 그 스택이 Artin 스택임을 증명하게
해 주는 변형 이론적 판정법, 이른바 Artin 판정법을 제시한다.
\end{quote}
\end{itemize}




\section{책과 온라인 강의록}
\label{guide-section-books}


\begin{itemize}
\item Laumon, Moret-Bailly: \emph{대수적 스택} \cite{LM-B}
\begin{quote}
이 책은 현재 스택에 관한 가장 포괄적인 참고문헌으로, 많은 기초 결과를
담고 있다. 독자가 대수공간에 익숙하다고 가정하며 Knutson의 책
\cite{Kn}을 자주 인용한다. 제12장에는 대수적 스택의 매끄러운-에탈
사이트의 함자성에 관한 오류가 있다. 이 오류는 Martin Olsson이 수정했으며
(\cite{olsson_sheaves} 참조), 나머지 장의 결과는 약간 수정해야 할 수는
있어도 올바르므로 걱정할 필요는 없다.
\end{quote}
\item 스택 프로젝트 저자들: \emph{스택 프로젝트}
\cite{stacks-project}.
\begin{quote}
바로 지금 읽고 있는 책이다!
\end{quote}
\item Anton Geraschenko:
\emph{Martin Olsson의 스택 강의를 위한 강의록} \cite{olsson_stacks}
\begin{quote}
이 강좌는 대수적 스택을 도입하기 전에 대수공간의 이론을 체계적으로
전개한다. 대수적 스택은 무려 제27강에서 처음 정의된다. 기본 성질뿐 아니라,
Deligne--Mumford 조건과 비분기 대각사상 조건의 동치, Artin 스택의
매끄러운-에탈 사이트, 준연접 층 이론, Keel--Mori 정리, 코호몰로지적
하강법, gerbe 및 Brauer 군과의 관계도 다룬다. 연습문제도 몇 가지 있다.
\end{quote}
\item
Behrend, Conrad, Edidin, Fantechi, Fulton, G\"ottsche, Kresch:
\emph{대수적 스택}, 현재 집필 중인 책의 온라인 강의록
\cite{stacks_book}
\begin{quote}
이 책의 목표는 고급 배경지식을 가정하지 않고 예제와 응용에 초점을 맞추어
스택을 친근하게 소개하는 것이다. \cite{LM-B}와 달리 독자가 대수공간
이론을 이미 소화했다고 가정하지 않는다. 대신 대수공간을 특별한 경우로
포함하도록 Deligne--Mumford 스택을 도입하며, 목표의 일부는 \cite{DM}의
주장들을 증명하기에 충분한 이론을 전개하는 것이다. Artin 스택의 일반
이론은 제2부에서 전개될 예정이다. 현재 Kresch의 웹사이트에는 책의 일부만
공개되어 있다.
\end{quote}
\item Olsson, Martin: \emph{대수공간과 스택}, \cite{olsson_book}
\begin{quote}
Hartshorne의 대수기하학 책을 익힌 독자의 수준에서 출발하는 대수공간과
대수적 스택 입문서로, 적극 추천한다.
\end{quote}
\end{itemize}


\section{스택의 기초에 관한 관련 참고문헌}
\label{guide-section-related}


\begin{itemize}
\item Vistoli:
\emph{Grothendieck 위상, 올범주 및 하강 이론에 관한 강의록}
\cite{vistoli_fga}
\begin{quote}
엄밀한 증명과 더불어 올범주, 스택, fpqc 위상에서의 하강 이론에 관한
유용한 사실들을 담고 있다.
\end{quote}
\item Knutson: \emph{대수공간} \cite{Kn}
\begin{quote}
Michael Artin의 지도 아래 쓴 박사학위 논문에서 발전한 이 책은 대수공간
이론의 기초를 담고 있다. \cite{LM-B}는 이 책을 자주 인용한다.
대수공간에 관한 Artin의 다음 논문들도 참조하라:
\cite{Artin-Algebraic-Approximation},
\cite{ArtinI}, \cite{Artin-Implicit-Function},
\cite{ArtinII}, \cite{Artin-Construction-Techniques},
\cite{Artin-Algebraic-Spaces}, \cite{Artin-Theorem-Representability},
그리고 \cite{ArtinVersal}.
\end{quote}
\item Grothendieck 외, \emph{토포스 이론과 스킴의 에탈 코호몰로지 I, II,
III}, SGA4로도 알려져 있다 \cite{SGA4}
\begin{quote}
제1권은 우주, 사이트, 올범주에 관한 많은 일반적 사실을 담고 있다.
“champ”(프랑스어로 “stack”)라는 단어는 Deligne의 Expos\'e XVIII에
등장한다.
\end{quote}
\item Jean Giraud: \emph{비가환 코호몰로지} \cite{giraud}
\begin{quote}
이 책은 일반 사이트 위의 올범주, 스택, 토서, gerbe를 논하지만 대수적
스택은 다루지 않는다. 예를 들어 $G$가 $X$ 위의 아벨 군의 층이면
$H^1(X, G)$를 $G$-토서와 동일시할 수 있는 것과 같은 방식으로
$H^2(X, G)$를 적절히 정의한 $G$-gerbe들의 집합과 동일시할 수 있다.
$G$가 아벨 군이 아닐 때에는 $H^2(X, G)$를 $G$-gerbe들의 집합으로
정의한다.
\end{quote}
\item Kelly와 Street: \emph{2-범주의 기본 요소에 대한 개관}
\cite{kelly-street}
\begin{quote}
스택의 범주는 2-범주를 이루며, 모든 2-사상이 가역인 단순한 유형의
2-범주이다. 이 논문은 일반적인 2-범주에 관한 참고문헌이다. 나는 이것을
사용해 본 적이 없어 얼마나 유용한지는 말할 수 없다. 또한
\cite{stacks-project}에도 2-범주의 기초가 일부 들어 있음을 유의하라.
\end{quote}
\end{itemize}




\section{문헌에 실린 논문}
\label{guide-section-papers}

\noindent
아래에는 스택과 대수공간에 관한 기본 결과를 담은 연구 논문을 열거한다.
요약문은 각 논문에서 스택 이론에 기여하는 결과만을 나타내려는 것이며,
많은 경우 이 결과들은 논문의 주된 목표에 부수적이다. 논문들을 여러
범주로 나누지만 일부 논문은 둘 이상의 범주에 속한다.


\subsection{변형 이론과 대수적 스택}
\label{guide-subsection-deformation-theory}

\noindent
Artin의 처음 세 논문은 스택 자체를 다루지는 않지만 강력한 결과를 담고
있으며, 그중 처음 두 논문은 \cite{ArtinVersal}에 필수적이다.
\begin{itemize}
\item Artin: \emph{완비 국소환 위 구조의 대수적 근사}
\cite{Artin-Algebraic-Approximation}
\begin{quote}
약한 가정 아래 모든 유효 형식 변형을 근사할 수 있음을 증명한다.
$F: (\Sch/S) \to (\textit{Sets})$가 국소 유한 표시인 반변 함자이고,
$S$가 체 위 유한형이거나 훌륭한 DVR이며, $s \in S$이고,
$\hat{\xi} \in F(\hat{\mathcal{O}}_{S, s})$가 유효 형식 변형이라고 하자.
그러면 모든 $n > 0$에 대하여 잉여적으로 자명한 에탈 근방
$(S', s') \to (S, s)$와 $\xi' \in F(S')$가 존재하여 $\xi'$와
$\hat{\xi}$가 $n$차까지 일치한다. 즉
$F(\mathcal{O}_{S, s} / \mathfrak m^n)$에서 같은 제한을 갖는다.
\end{quote}
\item
Artin: \emph{형식 모듈라이의 대수화 I} \cite{ArtinI}
\begin{quote}
약한 가정 아래 모든 유효 형식 버설 변형을 대수화할 수 있음을 증명한다.
$F: (\Sch/S) \to (\textit{Sets})$가 국소 유한 표시인 반변 함자이고,
$S$가 체 위 유한형이거나 훌륭한 DVR이며, $s \in S$가 국소 닫힌 점이라고
하자. 또한 $\hat A$가 완비 뇌터 국소 $\mathcal{O}_S$-대수이고,
그 잉여체 $k'$가 $k(s)$의 유한 확대이며,
$\hat{\xi} \in F(\hat A)$가 어떤 원소 $\xi_0 \in F(k')$의 유효 형식
버설 변형이라고 하자. 그러면 스킴 $X$가 존재하며 이는 $S$ 위 유한형이다.
또한 닫힌점 $x \in X$와 그 점에서 $k(x) = k'$인 잉여체, 그리고 원소
$\xi \in
F(X)$가 존재하여
$\hat{\mathcal{O}}_{X, x} \cong \hat{A}$인 동형이 있고, 이 동형 아래에서
$\xi$와 $\hat{\xi}$의 제한이 각 $F(\hat A /
\mathfrak m^n)$ 안에서 일치한다. $\hat{\xi}$가 보편
변형이면 대수화는 유일하다. Hilbert
스킴과 Picard 스킴의 표현가능성에 대한 응용도 제시한다.
\end{quote}
\item Artin: \emph{형식 모듈라이의 대수화 II}
\cite{ArtinII}
\begin{quote}
대략 말해, 닫힌 부분집합 $Y' \subset X'$를 $Y'$ 주위에서 형식적으로
국소 수축할 수 있으면, 전역 사상 $X' \to X$가 존재하여 $Y$를
수축하고 $X$는 대수공간임을 보인다.
\end{quote}
\item
Artin: \emph{버설 변형과 대수적 스택} \cite{ArtinVersal}
\begin{quote}
이 기념비적 논문은 \cite{Artin-Algebraic-Approximation}과
\cite{ArtinI}에서 이룬 연구를 바탕으로 한다. 변형 이론적 성질을
검증함으로써 스택의 대수성을 증명하게 해 주는 Artin 판정법을 도입한다.
좀 더 정확히 말하되 아주 정확하지는 않게 말하면, Artin은 극한을 보존하는
스택 $\mathcal{X}$의 점 $x \in \mathcal{X}(k)$ 주위에서 다음과 같이
표현을 구성한다. 스택 $\mathcal{X}$가 Schlessinger 판정법
(\cite{Sch})을 만족한다고 가정하면 형식 버설 변형
$\hat{\xi} \in \lim \mathcal{X}(\hat A / \mathfrak m^n)$이 존재하며
이는 $x$를 변형한다.
형식 변형이 유효하다고, 즉
$\mathcal{X}(\hat{A}) \to \lim \mathcal{X}(\hat A / \mathfrak m^n)$가
전단사라고 가정하면 유효 형식 버설 변형
$\xi \in \mathcal{X}(\hat A)$를 얻는다. \cite{ArtinI}의 결과를
사용하면 유한형 스킴 $U$와 원소 $\xi_U: U \to \mathcal{X}$를
얻으며, 이 원소는 한 점 $u \in U$에서 형식적으로 버설이고 그 점은
$x$ 위에 놓인다. 이제
$\mathcal{X}$가 몇 가지 조건, 즉 에탈 국소화 및 완비화와의 양립성과
구성 가능성 조건을 만족하는 변형 및 방해 이론을 갖는다고 가정하면,
제4절에서 형식 버설성이 열린 조건임을 보여 준다. 따라서 $U$를 줄이고
나면 $U
\to \mathcal{X}$는 매끄럽다.
또한 Artin은 스킴에 의한 fppf 표현을 허용하는 모든 스택이 스킴에 의한
매끄러운 표현을 허용함을 증명한다. 특히 평탄하고 분리이며 유한 표시인
군 스킴에 의한 몫 스택을 만들 수 있다.
\end{quote}
\item Conrad, de Jong: \emph{버설 변형의 근사}
\cite{conrad-dejong}
\begin{quote}
이 논문은 Popescu의 강력한 결과를 적용하여 Artin의 대수화 결과에
접근한다. $A$가 뇌터 환이고 $B$가 뇌터 $A$-대수이면, 사상
$A \to B$가 정칙 사상일 필요충분조건은 $B$가 매끄러운 $A$-대수들의
직접극한인 것이다. Popescu의 결과가 임의의 훌륭한 스킴 위에서 Artin의
근사 정리를 함의한다는 것은 어렵지 않게 알 수 있다. 훌륭함 가정은 국소환
$A$에 대하여 Hensel화에서 완비화로 가는 사상
$A^{\text{h}} \to \hat A$가 정칙임을 함의한다. 이 논문은 Popescu의
결과를 이용하여 \cite{ArtinI}의 주정리를 “군체”로 일반화하며, 이
일반화는 임의의 훌륭한 바탕 스킴과 임의의 점 $s \in S$에 대해 성립한다.
특히 \cite{ArtinVersal}의 결과는 임의의 훌륭한 바탕 위에서 성립한다.
저자들은 대수화의 에탈 국소 유일성과, 대상의 자기동형군이 대수화의
Hensel화에 자연스럽게 작용하는지도 논한다.
\end{quote}
\item Jason Starr: \emph{Artin 공리, 합성, 모듈라이 공간}
\cite{starr_artin}
\begin{quote}
이 논문은 대수화를 위한 Artin 공리들이 1-사상의 합성과 양립함을 확립한다.
\end{quote}
\item Martin Olsson: \emph{대수적 스택의 표현가능한 사상에 대한
변형 이론} \cite{olsson_deformation}
\begin{quote}
스킴의 사상에 관한 표준 변형 이론 결과를 여접 복합체를 사용하여 대수적
스택의 표현가능한 사상으로 일반화한다. 표현가능한 사상
$X \to \mathcal{X}$의 여접 복합체는 환 달린 토포스의 사상의 여접
복합체로 정의되지 않으므로, 이 결과들을 Illusie의 일반 이론의 귀결로
볼 수는 없다. 그 이유는 매끄러운-에탈 사이트가 함자적이지 않기 때문이다.
\end{quote}
\end{itemize}

\subsection{거친 모듈라이 공간}
\label{guide-subsection-coarse-moduli-spaces}

\noindent
거친 모듈라이 공간을 논하는 논문들이다.
\begin{itemize}
\item Keel, Mori: \emph{군체에서의 몫} \cite{K-M}
\begin{quote}
분리 Deligne--Mumford 스택이 거친 모듈라이 공간을 갖는다는 사실은
오랫동안 “민간전승”으로 알려져 있었던 듯하다. 여기서는 다음 정리의
엄밀하지만 간결한 증명을 제시한다. $\mathcal{X}$가 뇌터 바탕 스킴
위에서 국소 유한형인 Artin 스택이고 관성 스택
$I_\mathcal{X} \to \mathcal{X}$가 유한이면, 거친 모듈라이 공간
$\phi : \mathcal{X} \to Y$가 존재하여 $\phi$는 분리이고 $Y$는
$S$ 위 국소 유한형 대수공간이다. 관성이 유한하다는 가정은 정확히 올바른
조건이다. 거친 모듈라이 공간 $\phi : \mathcal{X} \to Y$가 존재하여
$\phi$가 분리일 필요충분조건은 관성이 유한한 것이다.
\end{quote}
\item Conrad: \emph{스택을 통한 Keel--Mori 정리} \cite{conrad}
\begin{quote}
Keel과 Mori의 논문 \cite{K-M}은 군체 언어로 쓰여 있어 이해하기
어렵다고 느끼는 이들이 있다. Brian Conrad는 정교한 스택 언어를
사용하면서도 매우 명료한, 스택 이론에 의한 증명을 제시한다. Conrad는
뇌터 가정도 제거한다.
\end{quote}
\item Rydh: \emph{유한군에 의한 몫과 거친 모듈라이 공간의 존재}
\cite{rydh_quotients}
\begin{quote}
Rydh는 $\mathcal{X}$가 어떤 바탕 위에서 유한 표시여야 한다는
\cite{K-M}과 \cite{conrad}의 가정을 제거한다.
\end{quote}
\item
Abramovich, Olsson, Vistoli: \emph{양의 표수에서의 순한 스택}
\cite{tame}
\begin{quote}
저자들은 유한 관성을 갖는 Artin 스택에서 거친 모듈라이 공간이
$\phi : \mathcal{X} \to Y$일 때, 준연접 층에 대한 $\phi_*$가
완전하면 이를 \emph{순한 Artin 스택}이라 정의한다. 유한 관성을 갖는
Artin 스택에 대하여 다음 조건들이 동치임을 증명한다. $\mathcal{X}$가
순하다. $\mathcal{X}$의 안정자들이 선형 환원적이다. 거친 모듈라이
공간 위에서 에탈 국소적으로 $\mathcal{X}$가 아핀 스킴을 선형 환원군
스킴으로 나눈 몫이다. 순한 Artin 스택에서는 거친 모듈라이 공간이 특히
좋은 성질을 갖는다. 예를 들어 거친 모듈라이 공간은 임의의 밑변환과
가환하지만, 유한 관성을 갖는 일반 Artin 스택의 거친 모듈라이 공간은
평탄 밑변환과만 가환한다.
\end{quote}
\item Alper: \emph{Artin 스택의 좋은 모듈라이 공간} \cite{alper_good}
\begin{quote}
무한 아핀 안정자군을 갖는 일반 Artin 스택은 필연적으로 비분리이며,
그 거친 모듈라이 공간은 흔히 존재하지 않는다. 가장 간단한 예는
$[\mathbf{A}^1 / \mathbf{G}_m]$이다. 여기서는 준콤팩트 사상
$\phi : \mathcal{X} \to Y$가 다음 두 조건을 만족할 때
\emph{좋은 모듈라이 공간}이라고 정의한다:
$\mathcal{O}_Y \to \phi_*
\mathcal{O}_\mathcal{X}$가 동형이고,
준연접 층에 대한 $\phi_*$가 완전하다.
이 개념은 \cite{tame}의 순한 Artin 스택을 일반화하는 동시에
Mumford의 기하 불변량 이론을 포괄한다. $G$가
$X \subset \mathbf{P}^n$에 선형으로 작용하는 환원군이면, 반안정
부분의 몫 스택에서 GIT 몫으로 가는 사상
$[X^{ss}/G] \to X//G$는 좋은 모듈라이 공간이다. 좋은 모듈라이 공간은
다음과 같은 좋은 기하적 성질을 많이 갖는다. (1) $\phi$는 전사이며
보편 닫힌 사상이고 보편 침수이다. (2) $\phi$는 $Y$의 점들을
$\mathcal{X}$의 점들과 폐포 동치까지 동일시한다. (3) $\phi$는
대수공간으로 가는 사상에 대해 보편적이다. (4) 좋은 모듈라이 공간은 임의의
밑변환 아래 안정적이다. (5) Artin 스택 위의 벡터다발이 좋은 모듈라이
공간으로 내려갈 필요충분조건은 닫힌점에서의 표현들이 자명한 것이다.
\end{quote}
\end{itemize}


\subsection{교차 이론}
\label{guide-subsection-intersection-theory}

\noindent
대수적 스택 위의 교차 이론을 논하는 논문들이다.
\begin{itemize}
\item
Vistoli: \emph{대수적 스택과 그 모듈라이 공간의 교차 이론}
\cite{vistoli_intersection}
\begin{quote}
이 논문은 Deligne--Mumford 스택에 대한 유리수 계수 교차 이론의 기초를
전개한다. $\mathcal{X}$가 분리 Deligne--Mumford 스택이면 유리수
계수 Chow 군 $\CH_*(\mathcal{X})$는 차원이 $k$인 정수적 닫힌
부분스택들로 생성되는 자유 아벨 군을 유리 동치로 나눈 것으로 정의한다.
평탄 당김, 고유 밀어내기, 정칙 국소 매장에 대한 일반화된 Gysin 준동형이
있다. $\phi : \mathcal{X} \to Y$가 모듈라이 공간, 즉 기하적 점에서
전단사인 고유 사상이면 유도된 밀어내기
$\CH_*(\mathcal{X}) \to
\CH_k(Y)$가 있으며 이는 동형이다.
\end{quote}
\item Edidin, Graham: \emph{등변 교차 이론}
\cite{edidin-graham}
\begin{quote}
이 논문의 목적은 몫 스택 $[X/G]$, 즉 대수군 $G$가 대수공간 $X$에
작용하여 얻는 몫에 대한 정수 계수 교차 이론, 다시 말해 $G$-등변
교차 이론을 $X$ 위에서 전개하는 것이다. 불변 순환만을 사용해 정의한
등변 Chow 군으로는
좋은 성질을 가진 이론이 나오지 않는다. 대신 $BG$의 경우에 대한 Totaro의
정의를 일반화하고, $V \to X$가 벡터다발이면 자연스럽게
$\CH_i(X) \cong \CH_i(V)$라는 사실에 착안하여 저자들은
$\CH_i^G(X)$를 다음과 같이 정의한다. $\dim(X) = n$이고
$\dim(G) = g$라 하자. 각 $i$에 대하여 $l$차원 $G$-표현 $V$를
고르되, $G$가 열린 부분집합 $U \subset V$에 자유롭게 작용하고 그
여집합의 여차원이 $d > n - i$가 되게 한다. 그러면
$X_G = [X \times U / G]$는 대수공간이며, 심지어 스킴으로 고를 수도
있다. 이제 $\CH_i^G(X) = \CH_{i + l - g}(X_G)$로 정의한다. 몫
스택에 대해서는
$\CH_i( [X/G]) = \CH_{i + g}^G(X) = \CH_{i + l}(X_G)$로 정의한다.
특히 $\CH_i([X/G]) = 0$인 것은 $i > \dim [X/G] = n - g$일
때이며,
$i \ll 0$일 때에는 0이 아닐 수 있다. 예를 들어
$\CH_i(B \mathbf{G}_m) = \mathbf{Z}$이며 이는 $i \le 0$에 대해
성립한다. 저자들은 이 등변 Chow
군들이 보통의 Chow 군과 같은 함자적 성질을 누림을 확립한다. 더 나아가
$[X / G] \cong [Y / H]$이면
$\CH_i([X/G]) = \CH_i([Y/H])$임을 보인다. 따라서 이 정의는 스택을
몫 스택으로 표현하는 방식에 의존하지 않는다.
\end{quote}
\item
Kresch: \emph{Artin 스택의 순환군} \cite{kresch_cycle}
\begin{quote}
Kresch는 임의의 Artin 스택에 대해 Chow 군을 정의한다. 몫 스택의
경우 이 정의는 \cite{edidin-graham}에서 Edidin과 Graham이 내린
정의와 일치한다. 아핀 안정자군을 갖는 대수적 스택에 대해서는 이 이론이
통상적인 성질들을 만족한다.
\end{quote}
\item Behrend와 Fantechi: \emph{내재적 법선뿔}
\cite{behrend-fantechi}
\begin{quote}
Behrend와 Fantechi는 Li와 Tian의 구성을 일반화하여
Deligne--Mumford 스택의 가상 기본류를 구성한다.
\end{quote}
\end{itemize}


\subsection{몫 스택}
\label{guide-subsection-quotient-stacks}

\noindent
몫 스택\footnote{문헌에서 \emph{몫 스택}은 흔히 $[X/G]$ 꼴의
스택을 뜻한다. 여기서 $X$는 대수공간이고 $G$는
$\text{GL}_n$의 부분군 스킴이며, 임의의 평탄 군 스킴을 허용하는 것은
아니다.}은 대수기하학자들이 연구하는 거의
모든 모듈라이 스택을 포함하는, Artin 스택의 매우 중요한 부분류를 이룬다.
몫 스택 $[X/G]$의 기하는 $G$-등변 기하, 곧 $X$의 등변 기하이다.
어떤 성질은 몫
스택에 대해 증명하기가 더 쉽고, 어떤 결과는 몫 스택에 대해서만 알려져
있다. 다음 논문들은 이런 질문을 다룬다. 대수적 스택은 언제 전역 몫
스택인가? 대수적 스택은 “국소적으로” 몫 스택인가?
\begin{itemize}
\item Laumon, Moret-Bailly: \cite[Chapter 6]{LM-B}
\begin{quote}
제6장은 대수적 스택의 국소 및 전역 구조에 관한 여러 사실을 담고 있다.
대수적 스택 $\mathcal{X}$가 $S$ 위에서 몫 스택 $[Y/G]$일
필요충분조건을 제시한다. 여기서 $Y$는 대수공간(각각 스킴,
아핀 스킴)이고 $G$는 유한군이며, 그 필요충분조건은
대수공간(각각 스킴, 아핀 스킴) $Y'$과 유한 에탈 사상
$Y'
\to \mathcal{X}$가 존재하는 것이다. $S$ 위의 임의의
Deligne--Mumford 스택과 $x : \Spec(K) \to \mathcal{X}$에 대해,
표현가능하고 에탈이며 분리인 사상
$\phi : [X/G] \to \mathcal{X}$가 존재한다. 여기서 $G$는 유한군이고
$S$ 위의 아핀 스킴에 작용하며,
$\Spec(K) = [X/G] \times_\mathcal{X} \Spec(K)$임을 보인다.
기하적으로 연결된 올을 갖는 표현의 존재도 자세히 논한다.
\end{quote}
\item
Edidin, Hassett, Kresch, Vistoli:
\emph{Brauer 군과 몫 스택} \cite{ehkv}
\begin{quote}
먼저 주어진 대수적 스택이 몫 스택이 되는 때에 관한 몇 가지 기본적이지만
그리 어렵지는 않은 사실을 확립한다. 이 논문에서는 대수적 스택이 항상
뇌터 스킴 위 유한형이라고 가정한다. 대수적 스택 $\mathcal{X}$에
대하여 다음 조건들은 동치이다. $\mathcal{X}$가 몫 스택이다. 모든
기하적 점에서 안정자가 올에 충실하게 작용하는 벡터다발
$V \to \mathcal{X}$가 존재한다. 벡터다발 $V \to \mathcal{X}$와
국소 닫힌 부분스택 $V^0 \subset V$가 존재하여, $V^0$가 표현가능하며
$\mathcal{X}$ 위로 전사한다. 대수공간에 의한 유한 평탄 덮개가 존재하면
대수적 스택이 몫 스택임을 확립한다. 일반점에서 안정자가 자명한 임의의
매끄러운 Deligne--Mumford 스택은 몫 스택이다.
$\mathbf{G}_m$-gerbe가 뇌터 스킴 $X$ 위에 있고
$\beta \in H^2(X, \mathbf{G}_m)$에 대응할 때, 그것이 몫 스택일
필요충분조건은 $\beta$가 Brauer 사상
$\text{Br}(X) \to \text{Br}'(X)$의 상에 속하는 것임을 보인다.
이를 이용하여 몫 스택이 아닌 비분리 Deligne--Mumford 스택을 만든다.
\end{quote}
\item Totaro: \emph{스킴과 스택의 해소 성질}
\cite{totaro_resolution}
\begin{quote}
모든 연접 층이 어떤 벡터다발의 몫이면 스택이 해소 성질을 갖는다고 한다.
첫 번째 주정리는 닫힌점에서 아핀 안정자군을 갖는 정규 뇌터 대수적
스택 $\mathcal{X}$에 대해 다음 두 조건이 동치라는 것이다:
(1) $\mathcal{X}$가 해소 성질을 갖는다. (2)
$\mathcal{X} = [Y/\text{GL}_n]$이고 $Y$는 준아핀이다.
$\mathcal{X}$가 체 위 유한형이면 (1), (2)는 다음 조건과도 동치이다:
(3) $\mathcal{X} = [\Spec(A)/G]$이고 $G$는 $k$ 위 유한형 아핀
군 스킴이다. 몫 스택이 해소 성질을 갖는다는 함의는 Thomason이 증명했다.
두 번째 주정리는 다음과 같다. $\mathcal{X}$가 체 위의 매끄러운
Deligne--Mumford 스택이고 유한하며 일반점에서 자명한 안정자군
$I_\mathcal{X}
\to \mathcal{X}$를 가지며, 그 거친 모듈라이 공간이
아핀 대각사상을 갖는 스킴이면 $\mathcal{X}$는 해소 성질을 갖는다.
또 다른 흥미로운 결과는 다음과 같다. 해소 성질을 만족하는 뇌터 대수적
스택 $\mathcal{X}$가 주어지면, $\mathcal{X}$가 아핀 대각사상을
가질 필요충분조건은 닫힌점들이 아핀 안정자를 갖는다는 것이다.
\end{quote}
\item Kresch: \emph{Deligne--Mumford 스택의 기하}
\cite{kresch_geometry}
\begin{quote}
이 논문은 체 위 유한형 Deligne--Mumford 스택의 일반 구조 정리를
요약하고 몫 스택에 관한 몇 가지 흥미로운 결과를 담고 있다.
준사영 거친 모듈라이 공간을 갖는 매끄럽고 분리이며 일반점에서 순한
Deligne--Mumford 스택은 몫 스택 $[Y/G]$임을 보인다. 여기서 $Y$는
준사영이고 $G$는 대수군이다. 거친 모듈라이 공간이 스킴인
Deligne--Mumford 스택 $\mathcal{X}$에 대해, $\mathcal{X}$가
Zariski 국소적으로 몫 스택일 필요충분조건은 스킴을 유한군으로 나눈 스택 몫들로 이루어진
Zariski 열린 덮개를 갖는 것이다. $\mathcal{X}$가 표수 0인 체 위에서
고유인 Deligne--Mumford 스택이고 거친 모듈라이 공간이 $Y$이면 다음
조건들은 동치이다. $Y$가 사영이고 $\mathcal{X}$가 몫 스택이다.
$Y$가 사영이고 $\mathcal{X}$가 생성 층을 갖는다. $\mathcal{X}$가
사영 거친 모듈라이 공간을 갖는 매끄럽고 고유인 DM 스택으로 닫힌 매장을
허용한다. 이에 착안하여, 사영 거친 모듈라이 공간을 갖는 매끄럽고 고유인
Deligne--Mumford 스택으로 닫힌 매장이 존재할 때 해당
Deligne--Mumford 스택이 \emph{사영적}이라고 정의한다.
\end{quote}
\item Kresch, Vistoli \emph{Deligne--Mumford 스택의 덮개와
Brauer 사상의 전사성} \cite{kresch-vistoli}
\begin{quote}
표수 0이고 $n$이 고정되었을 때 다음 두 명제가 동치임을 보인다:
(1) 모든 $n$차원 매끄러운 Deligne--Mumford 스택은 몫 스택이다.
(2) $n$차원 매끄러운 스킴에서 Azumaya Brauer 군과 코호몰로지적
Brauer 군이 일치한다.
\end{quote}
\item Kresch: \emph{Artin 스택의 순환군} \cite{kresch_cycle}
\begin{quote}
체 위 유한형인 축약 Artin 스택이 아핀 안정자군을 가지면 몫 스택들에
의한 층화를 갖는다는 것을 보인다.
\end{quote}
\item Abramovich-Vistoli:
\emph{안정 사상 공간의 콤팩트화} \cite{abramovich-vistoli}
\begin{quote}
보조정리 2.2.3은 임의의 분리 Deligne--Mumford 스택이 거친 모듈라이
공간 위에서 에탈 국소적으로 몫 스택 $[U/G]$임을 확립한다. 여기서
$U$는 아핀이고 $G$는 유한군이다. \cite[Theorem 2.12]{olsson_homstacks}에
따르면 이 논증에서 $G$는 실제로 안정자군이다.
\end{quote}
\item Abramovich, Olsson, Vistoli:
\emph{양의 표수에서의 순한 스택} \cite{tame}
\begin{quote}
이 논문은 순한 Artin 스택이 거친 모듈라이 공간 위에서 에탈 국소적으로
아핀 스킴을 안정자군으로 나눈 몫 스택임을 보인다.
\end{quote}
\item Alper: \emph{Artin 스택의 국소 몫 구조에 관하여}
\cite{alper_quotient}
\begin{quote}
Artin 스택 $\mathcal{X}$와 선형 환원적 안정자를 갖는 닫힌점
$x
\in \mathcal{X}$가 주어지면, 에탈 사상
$[V/G_x]
\to \mathcal{X}$가 존재하며 $V$는 대수공간일 것이라고
추측한다. 이 추측을
뒷받침하는 몇 가지 증거를 제시한다. \cite{tame}의 발상에 기초한
간단한 변형 이론 논증은 이 명제가 형식 국소적으로 참임을 보인다.
Luna의 에탈 절편 정리에 대한 스택 이론적 증명도 제시한다. 즉 스택
$\mathcal{X} = [\Spec(A)/G]$에서 $G$가 선형 환원적이면, GIT 몫
$\Spec(A^G)$ 위에서 에탈 국소적으로 $\mathcal{X}$는 안정자에 의한
몫 스택임을 증명한다.
\end{quote}
\end{itemize}

\subsection{코호몰로지}
\label{guide-subsection-cohomology}

\noindent
대수적 스택 위의 층의 코호몰로지를 논하는 논문들이다.
\begin{itemize}
\item Olsson: \emph{Artin 스택 위의 층} \cite{olsson_sheaves}
\begin{quote}
이 논문은 준연접 층과 구성 가능 층의 이론을 전개하고 기본적인
코호몰로지적 성질을 증명한다. \cite{LM-B}에서 매끄러운-에탈 사이트의
함자성에 관해 범한 오류를 바로잡는다. 여접 복합체도 구성한다. 또한 다음
정리들을 증명한다. 고유 사상에 대한 Grothendieck 기본 정리,
Grothendieck 존재 정리, Zariski 연결성 정리, 그리고 연접 층과 구성
가능 층의 고유 밀어내기에 대한 유한성 정리.
\end{quote}
\item Behrend: \emph{대수적 스택의 유도 $l$-진 범주}
\cite{behrend_derived}
\begin{quote}
대수적 스택에 대한 Lefschetz 대각합 공식을 증명한다.
\end{quote}
\item Behrend: \emph{스택의 코호몰로지} \cite{behrend_cohomology}
\begin{quote}
미분 가능 스택의 de Rham 코호몰로지와 위상 스택의 특이 코호몰로지를
정의한다.
\end{quote}
\item Faltings:
\emph{고유 fppf 스택의 연접 코호몰로지의 유한성}
\cite{faltings_finiteness}
\begin{quote}
고유 사상에 따른 연접 층의 직접상이 연접임을 증명한다.
\end{quote}
\item Abramovich, Corti, Vistoli:
\emph{뒤틀린 다발과 허용 가능 덮개} \cite{acv}
\begin{quote}
부록에는 순한 Deligne--Mumford 스택의 에탈 코호몰로지에 대한
고유 밑변환 정리가 들어 있다.
\end{quote}
\end{itemize}


\subsection{스킴에 의한 유한 덮개의 존재}
\label{guide-subsection-finite-covers}

\noindent
스킴에 의한 Deligne--Mumford 스택의 유한 덮개가 존재한다는 것은 중요한
결과이다. Deligne--Mumford 스택의 교차 이론에서는 표현가능하지 않은
사상의 고유 밀어내기를 정의하는 데 필수적으로 쓰인다.
$\overline{\mathcal{M}}_g$에 관한 여러 결과는 Looijenga가 증명한
\emph{매끄러운} 스킴에 의한 유한 덮개의 존재에 의존한다. 이 방향의
최초 결과는 등변 상황을 다루는
\cite[Theorem 6.1]{seshadri_quotients}일 수도 있다.
\begin{itemize}
\item Vistoli: \emph{대수적 스택과 그 모듈라이 공간의 교차 이론}
\cite{vistoli_intersection}
\begin{quote}
$\mathcal{X}$가 모듈라이 공간, 즉 기하적 점에서 전단사인 고유 사상을
갖는 Deligne--Mumford 스택이면 유한 사상 $X \to \mathcal{X}$가
존재하며 이는 어떤 스킴 $X$에서 온다.
\end{quote}
\item Laumon, Moret-Bailly: \cite[Chapter 16]{LM-B}
\begin{quote}
Zariski 주정리의 응용으로 정리 16.6은 다음을 확립한다.
$\mathcal{X}$가 뇌터 스킴 위 유한형 Deligne--Mumford 스택이면,
유한하고 전사이며 일반점에서 에탈인 사상 $Z \to
\mathcal{X}$가
존재하고 $Z$는 스킴이다. 따름정리 16.6.2에서는 임의의 뇌터 정규
대수공간이 대수공간 몫 $X'/G$와 동형임도 보인다. 여기서 유한군 $G$는
정규 스킴 $X$에 작용한다.
\end{quote}
\item Edidin, Hassett, Kresch, Vistoli:
\emph{Brauer 군과 몫 스택} \cite{ehkv}
\begin{quote}
정리 2.7의 내용은 다음과 같다. $\mathcal{X}$가 뇌터 바탕 스킴
$S$ 위 유한형 대수적 스택이면, 대각사상
$\mathcal{X} \to \mathcal{X} \times_S \mathcal{X}$가 준유한일
필요충분조건은 유한 전사 사상 $X \to F$가 존재하고 그 사상이 어떤
스킴 $X$에서 오는 것이다.
\end{quote}
\item Kresch, Vistoli: \emph{Deligne--Mumford 스택의 덮개와
Brauer 사상의 전사성} \cite{kresch-vistoli}
\begin{quote}
준사영 거친 모듈라이 공간을 갖고 체 위 유한형인 임의의 매끄러운 분리
Deligne--Mumford 스택이 매끄러운 준사영 스킴에 의한 유한 평탄 덮개를
갖는다는 것을 증명한다.
\end{quote}
\item Olsson: \emph{Artin 스택의 고유 덮개에 관하여}
\cite{olsson_proper}
\begin{quote}
$\mathcal{X}$가 $S$ 위 분리 유한형 Artin 스택이면, 고유 전사 사상
$X \to \mathcal{X}$가 존재함을 증명한다. 이 사상은 어떤 스킴 $X$에서
오며 그 스킴은 $S$ 위 준사영이다. 응용으로 Olsson은 고유 사상에
따른 직접상 층의 연접성과 구성 가능성을 증명한다. 또 다른 응용으로 고유
Artin 스택에 대한 Grothendieck 존재 정리를 증명한다.
\end{quote}
\item Rydh: \emph{대수공간과 스택의 뇌터 근사}
\cite{rydh_approx}
\begin{quote}
이 논문의 정리 B는 다음과 같다. $X$가 준유한이고 분리인 대각사상을 갖는
준콤팩트 대수적 스택이라고 하자(각각, 준콤팩트이고 분리인 대각사상을
갖는 준콤팩트 Deligne--Mumford 스택이라고 하자). 그러면
스킴 $Z$와 유한하고 유한 표시이며 전사인 사상 $Z \to X$가 존재하여,
어떤 조밀한 준콤팩트 열린 부분스택 $U \subset X$ 위에서 이 사상은
평탄하다(각각 에탈이다).
\end{quote}
\end{itemize}


\subsection{경직화}
\label{guide-subsection-rigidification}

\noindent
경직화는 관성에서 평탄 부분군을 제거하는 과정이다. 예를 들어 $X$가
사영 다양체이면 Picard 스택에서 Picard 스킴으로 가는 사상은
자기동형군 $\mathbf{G}_m$의 경직화이다.
\begin{itemize}
\item Abramovich, Corti, Vistoli:
\emph{뒤틀린 다발과 허용 가능 덮개} \cite{acv}
\begin{quote}
$\mathcal{X}$가 $S$ 위의 대수적 스택이고 $H$가 $S$ 위의 평탄하고
유한 표시이며 분리인 군 스킴이라 하자. 모든 대상
$\xi \in \mathcal{X}(T)$에 대하여 매장
$H(T) \hookrightarrow \text{Aut}_{\mathcal{X}(T)}(\xi)$이 존재하며
다음 의미에서 당김과 양립한다고 가정하자. 모든 화살
$\phi : \xi \rightarrow \xi'$가 $f: T \rightarrow T'$ 위에 놓이고,
모든 $g \in H(T')$에
대하여 $g \circ \phi = \phi \circ f^*g$이다. 그러면 대수적 스택
$\mathcal{X}/H$와 fppf gerbe인 사상
$\rho : \mathcal{X} \rightarrow \mathcal{X}/H$가 존재하며,
모든 $\xi \in \mathcal{X}(T)$에 대해 사상
$\text{Aut}_{\mathcal{X}(T)} (\xi)
\rightarrow \text{Aut}_{\mathcal{X}/H (T)} (\xi) $
은 전사이고 그 핵은 $H(T)$이다.
\end{quote}
\item Romagny: \emph{스택 위의 군 작용과 응용}
\cite{romagny_actions}
\begin{quote}
군 작용이 경직화와 관련하여 어떻게 행동하는지 논한다.
\end{quote}
\item Abramovich, Graber, Vistoli:
\emph{Deligne--Mumford 스택의 Gromov--Witten 이론} \cite{agv}
\begin{quote}
부록에서는 \cite{acv}에서와 같은 경직화를 두 가지 대안적 해석과 함께
요약한다. 이 논문은 닫힌 부분스택을 따라 대수적 스택을 붙이는 구성과
선다발의 근을 취하는 구성도 담고 있다.
\end{quote}
\item
Abramovich, Olsson, Vistoli: \emph{양의 표수에서의 순한 스택}
(\cite{tame})
\begin{quote}
부록은 관성의 평탄 부분군 스택 $H \subset I_\mathcal{X}$가
정규이지만 반드시 중심적이지는 않은 더 복잡한 상황을 다룬다.
\end{quote}
\end{itemize}

\subsection{스택 곡선}
\label{guide-subsection-stacky-curves}

\noindent
스택 곡선을 논하는 논문들이다.
\begin{itemize}
\item Abramovich, Vistoli: \emph{안정 사상 공간의 콤팩트화}
\cite{abramovich-vistoli}
\begin{quote}
이 논문은 \emph{뒤틀린 곡선}을 도입한다. 안정 곡선에서 대수적 스택으로
가는 안정 사상들의 모듈라이 공간은 일반적으로 콤팩트하지 않다. 저자들은
뒤틀린 곡선에서 오는 사상을 이용하여, 표적이 사영 거친 모듈라이 공간을
갖는 순한 Deligne--Mumford 스택일 때 고유인 모듈라이 스택을 구성한다.
\end{quote}
\item Behrend, Noohi: \emph{Deligne--Mumford 곡선의 균일화}
\cite{behrend-noohi}
\begin{quote}
Deligne--Mumford 해석 곡선의 균일화 정리를 증명한다.
\end{quote}
\end{itemize}

\subsection{Hilbert, Quot, Hom 및 분지다양체 스택}
\label{guide-subsection-hilbert-quot-hom}

\noindent
Hilbert 스킴과 그와 비슷한 대상을 논하는 논문들이다.
\begin{itemize}
\item Vistoli: \emph{Hilbert 스택과 족의 모듈라이 이론}
\cite{vistoli_hilbert}
\begin{quote}
$\mathcal{X}$가 국소 뇌터이며 국소 분리인 대수공간 $S$ 위에서
분리이고 국소 유한형인 대수적 스택이면, Vistoli는 고유 스킴에서 오는
유한 비분기 사상을 매개화하는 Hilbert 스택
$\mathcal{H}\text{ilb}(\mathcal{F} / S)$를 정의한다.
$\mathcal{H}\text{ilb}(\mathcal{F} / S)$가 대수적 스택이라는
주장은 증명 없이 제시된다. 그 귀결로, 위와 같은 $\mathcal{X}$에 대해
Hom 스택 $\mathcal{H} \text{om}_S(T, \mathcal{X})$가 대수적
스택임을 증명한다. 여기서 $T$는 $S$ 위에서 고유이고 평탄하다.
\end{quote}
\item Olsson, Starr: \emph{Deligne--Mumford 스택의 Quot 함자}
\cite{olsson-starr}
\begin{quote}
$\mathcal{X}$가 대수공간 $S$ 위에서 분리이고 국소 유한 표시인
Deligne--Mumford 스택이며 $\mathcal{F}$가 국소 유한 표시
$\mathcal{O}_\mathcal{X}$-가군이면, Quot 함자
$\text{Quot}(\mathcal{F} / \mathcal{X} / S)$는 $S$ 위에서 분리이고
국소 유한 표시인 대수공간으로 표현된다. 이 논문은 생성 층도 정의하며,
유한군으로 스킴을 나눈 전역 몫 스택인 순하고 분리인
Deligne--Mumford 스택에 대해 생성 층의 존재를 증명한다.
\end{quote}
\item Olsson: \emph{Hom-스택과 스칼라 제한}
\cite{olsson_homstacks}
\begin{quote}
$\mathcal{X}$와 $\mathcal{Y}$가 유한 대각사상을 가지며 대수공간
$S$ 위에서 국소 유한 표시인 Artin 스택이라고 하자. $\mathcal{X}$는
$S$ 위에서 고유이고 평탄하며, $S$ 위에서 fppf 국소적으로
$\mathcal{X}$는 대수공간에 의한 유한하고 유한 표시이며 평탄인 덮개를
허용한다고 하자. 예를 들어
$\mathcal{X}$가 Deligne--Mumford 스택이거나 순한 Artin 스택이면
된다. 그러면 $\Hom_S(\mathcal{X}, \mathcal{Y})$는 $S$ 위에서
국소 유한 표시인 Artin 스택이다.
\end{quote}
\item Alexeev와 Knutson: \emph{분지다양체의 완비 모듈라이 공간}
(\cite{alexeev-knutson})
\begin{quote}
저자들은 $\mathbf{P}^n$의 분지다양체를 유한 사상
$X \rightarrow \mathbf{P}^n$으로 정의하며, 그 원천 $X$는
\emph{축약} 스킴이다. Hilbert
다항식과 $i$차원 성분들의 전체 차수를 고정한 분지다양체들의 모듈라이
스택이 유한 안정자를 갖는 고유 Artin 스택임을 증명한다. 분지다양체
스택을 Hilbert 스킴, Chow 스킴, 안정 사상의 모듈라이 공간과 비교한다.
\end{quote}
\item Lieblich: \emph{연접 대수 스택에 관한 논평}
\cite{lieblich_remarks}
\begin{quote}
이 논문은 스킴 $Y$ 위에서 Alexeev와 Knutson의 분지다양체 스택을
일반화한다. 이를 위해 $Y$의 구조층 위 대수들의 스택으로 이 스택을
구성한다. $\text{Quot}$ 공간과 $\Hom$ 공간의 존재 증명도 제시한다.
\end{quote}
\item Starr: \emph{Artin 공리, 합성, 모듈라이 공간}
\cite{starr_artin}
\begin{quote}
주요 결과의 응용으로 Vistoli의 Hilbert 스택
\cite{vistoli_hilbert}과 Alexeev--Knutson의 분지다양체 스택
\cite{alexeev-knutson}을 동시에 일반화한다. $\mathcal{X}$가 훌륭한
스킴 $S$ 위에서 국소 유한형이며 유한 대각사상을 갖는 대수적 스택이면,
스택 $\mathcal{H}$는 사상 $g: T \rightarrow \mathcal{X}$를
매개화한다. 여기서 $T$는 $G$-풍부 선다발 $L$을 갖는 고유
대수공간이다. 이 스택은 $S$ 위 국소 유한형 Artin 스택이다.
\end{quote}
\item Lundkvist와 Skjelnes:
\emph{Grothendieck의 Hilbert 함자의 비유효 변형}
\cite{lundkvist-skjelnes}
\begin{quote}
비유효 변형이 존재하므로 비분리 스킴의 Hilbert 함자는 표현되지 않음을
보인다.
\end{quote}
\item Halpern-Leistner와 Preygel:
\emph{사상 스택과 고유성의 범주론적 개념}
\cite{HL-P}
\begin{quote}
Olsson의 논문에서 쓰인 가정보다 더 일반적이거나 적어도 다른, 원천과
표적에 관한 몇 가지 가정 아래 Hom 스택이 대수적임을 증명한다.
\end{quote}
\end{itemize}



\subsection{토릭 스택}
\label{guide-subsection-toric}

\noindent
토릭 다양체가 스킴 이론에서 예와 반례를 제공하는 것과 마찬가지로,
토릭 스택의 기하와 그 스택 부채의 조합론 사이에는 사전과 같은 대응이
있다. 따라서 토릭 스택은 훌륭한 예들의 부류이자 추측을 시험하는 자연스러운
무대이다.
\begin{itemize}
\item Borisov, Chen, Smith: \emph{토릭 Deligne--Mumford 스택의
오비폴드 Chow 환} \cite{bcs}
\begin{quote}
토릭 다양체에 대한 Cox의 구성에서 영감을 얻어, 이 논문은
\emph{스택 부채}라 불리는 조합론적 대상에 대응하는 명시적 몫 스택으로
매끄러운 토릭 DM 스택을 정의한다.
\end{quote}
\item Iwanari: \emph{토릭 스택의 범주} \cite{iwanari_toric}
\begin{quote}
이 논문은 \emph{토릭 삼중항}을 다음 자료로 정의한다. 매끄러운
Deligne--Mumford 스택 $\mathcal{X}$와 조밀한 상을 갖는 열린 매장
$\mathbf{G}_m \hookrightarrow \mathcal{X}$가 주어지며, 따라서
$\mathcal{X}$는 오비폴드이다. 또한 작용
$\mathcal{X} \times \mathbf{G}_m \rightarrow \mathcal{X}$가
주어진다. 토릭 삼중항의 2-범주와 스택
부채의 1-범주가 동치임을 보인다. 토릭 삼중항과 \cite{bcs}의 매끄러운
토릭 DM 스택 정의 사이의 관계도 논한다.
\end{quote}
\item Iwanari: \emph{토릭 스택의 정수 Chow 환}
\cite{iwanari_chow}
\begin{quote}
토릭 다양체에 대한 Cox의 $\Delta$-모음을 토릭 오비폴드로 일반화한다.
\end{quote}
\item Perroni: \emph{토릭 Deligne--Mumford 스택에 관한 짧은 글}
\cite{perroni}
\begin{quote}
Cox의 $\Delta$-모음과 Iwanari의 논문 \cite{iwanari_chow}을 일반적인
매끄러운 토릭 DM 스택으로 일반화한다.
\end{quote}
\item Fantechi, Mann, Nironi: \emph{매끄러운 토릭 DM 스택}
\cite{fmn}
\begin{quote}
이 논문은 매끄러운 토릭 DM 스택을 다음과 같이 정의한다. 매끄러운 DM
스택 $\mathcal{X}$에 DM 토러스 $\mathcal{T}$가 작용한다. 여기서
DM 토러스란 $T \times BG$와 동형인 Picard 스택을 뜻하며 $G$는
유한이다. 이 작용은 $\mathcal{T}$와 동형인 열린 조밀 궤도를 갖는다.
저자들은 “아래에서 위로 올라가는 서술”을 제시하고 매끄러운 토릭 DM
스택과 스택 부채 사이의 동치를 증명한다.
\end{quote}
\item Geraschenko와 Satriano: \emph{토릭 스택 I 및 II}
\cite{gs_toric1} 및 \cite{gs_toric2}
\begin{quote}
이 논문들은 토릭 스택을 토릭 다양체를 그 토러스의 부분군으로 나눈
스택 몫으로 정의한다. 일반점에서 스택적인 토릭 스택은 토릭 스택의 토러스
불변 부분스택으로 정의한다. 이 정의는 앞선 토릭 스택 정의들을 포괄하고
확장한다. 첫 번째 논문은 스택 부채의 조합론과 대응하는 스택의 기하 사이의
사전을 전개한다. 또한 \cite{perroni}의 해석을 일반화하여 매끄러운 토릭
스택의 모듈라이 해석을 제시한다. 두 번째 논문은 토릭 스택의 내재적
특성화를 증명한다.
\end{quote}
\end{itemize}

\subsection{형식 함수 정리와 Grothendieck 존재 정리}
\label{guide-subsection-theorem-formal-functions}

\noindent
이 논문들은 형식 함수 정리 \cite[III.4.1.5]{EGA}와 Grothendieck 존재
정리 \cite[III.5.1.4]{EGA}를 일반화한다. 전자는 때때로 고유 사상에
대한 Grothendieck 기본 정리라고도 불린다.
\begin{itemize}
\item Knutson: \emph{대수공간} \cite[Chapter V]{Kn}
\begin{quote}
이 정리들을 대수공간으로 일반화한다.
\end{quote}
\item Abramovich-Vistoli: \emph{안정 사상 공간의 콤팩트화}
\cite[A.1.1]{abramovich-vistoli}
\begin{quote}
이 정리들을 순한 Deligne--Mumford 스택으로 일반화한다.
\end{quote}
\item Olsson과 Starr: \emph{Deligne--Mumford 스택의 Quot 함자}
\cite{olsson-starr}
\begin{quote}
이 정리들을 분리 Deligne--Mumford 스택으로 일반화한다.
\end{quote}
\item Olsson: \emph{Artin 스택의 고유 덮개에 관하여}
\cite{olsson_proper}
\begin{quote}
고유 Artin 스택으로 일반화한다.
\end{quote}
\item Conrad: \emph{Artin 스택 위의 형식 GAGA} \cite{conrad_gaga}
\begin{quote}
고유 Artin 스택으로 일반화하고 형식 GAGA 정리를 증명한다.
\end{quote}
\item Olsson: \emph{Artin 스택 위의 층} \cite{olsson_sheaves}
\begin{quote}
고유 Artin 스택으로의 일반화에 대한 또 다른 증명을 제시한다.
\end{quote}
\end{itemize}


\subsection{스택 위의 군 작용}
\label{guide-subsection-group-actions}

\noindent
대수적 스택 위의 군 작용은 자연스럽게 나타난다. 예를 들어 대칭군
$S_n$은 $\overline{\mathcal{M}}_{g, n}$에 작용하고, 군 $G$가 스킴
$X$에 작용하면 $G$의 $\text{Aut}(X)$ 안의 정규화군이 $[X/G]$에
작용한다. 또한 스택 위의 토러스 작용은 Gromov--Witten 이론에 자주
등장한다.
\begin{itemize}
\item Romagny: \emph{스택 위의 군 작용과 응용}
\cite{romagny_actions}
\begin{quote}
이 논문은 군이 대수적 스택에 작용한다는 말의 의미를 정밀하게 만들고,
대수적 스택 위 군 스킴의 작용에 대해 고정점의 존재와 몫의 존재를
증명한다. Romagny의 이전 강의록 \cite{romagny_notes}도 참조하라.
\end{quote}
\end{itemize}

\subsection{선다발의 근 취하기}
\label{guide-subsection-root-stacks}

\noindent
이 유용한 구성은 Cadman과 Abramovich, Graber, Vistoli가 서로 독립적으로
발견했다. 스킴 $X$와 그 위의 유효 Cartier 인자 $D$가 주어지면,
$r$번째 근 스택은 $X$ 위에서 $D$를 따라 분기하고
$\mu_r$ 안정자를 $D$ 위에 가지며 $D$ 밖에서는 스킴과 같은 Artin
스택이다.
\begin{itemize}
\item Charles Cadman
\emph{스택을 이용해 곡선에 접촉 조건을 부과하기}
\cite{cadman}
\item Abramovich, Graber, Vistoli: \emph{Deligne--Mumford 스택의
Gromov--Witten 이론} \cite{agv}
\end{itemize}

\subsection{그 밖의 논문}
\label{guide-subsection-other}

\noindent
그 밖의 여러 논문이다.
\begin{itemize}
\item Lieblich: \emph{뒤틀린 층의 모듈라이} \cite{lieblich_twisted}
\begin{quote}
이 논문은 gerbe와 뒤틀린 층을 요약한다.
$\mathcal{X} \rightarrow X$가 $\mu_n$-gerbe이고 $X$가 매끄럽고
연결된 기하적 올을 갖는 사영 상대 곡면이면, 반안정
$\mathcal{X}$-뒤틀린 층의 스택이 $S$ 위에서 국소 유한 표시인 Artin
스택임을 보인다. 또한 Artin 스택 위 층의 연관점과 순수성 이론도 전개한다.
\end{quote}
\item Lieblich, Osserman:
\emph{스택의 함자적 복원 정리}
\cite{lieblich-osserman}
\begin{quote}
대수적 스택을 그에 대응하는 함자로부터 언제 복원할 수 있는지에 관한
놀랍고 흥미로운 결과들을 증명한다.
\end{quote}
\item David Rydh:
\emph{대수공간과 스택의 뇌터 근사}
\cite{rydh_approx}
\begin{quote}
준유한 대각사상을 갖는 모든 준콤팩트 대수적 스택을 뇌터 스택으로
근사할 수 있음을 보인다. Chevalley, Serre, Zariski, Chow의 결과에서
뇌터 가정을 제거하는 데 응용된다.
\end{quote}
\end{itemize}



\section{다른 분야의 스택}
\label{guide-section-stacks-areas}

\begin{itemize}
\item Behrend와 Noohi: \emph{Deligne--Mumford 곡선의 균일화}
\cite{behrend-noohi}
\begin{quote}
위상 스택, 해석 스택, 대수적 스택을 개관하고 비교한다.
\end{quote}
\item Behrang Noohi: \emph{위상 스택의 기초 I} \cite{noohi}
\item David Metzler: \emph{위상 스택과 매끄러운 스택} \cite{metzler}
\end{itemize}


\section{고차 스택}
\label{guide-section-higher-stacks}

\begin{itemize}
\item Lurie: \emph{고차 토포스 이론} \cite{lurie_topos}
\item Lurie: \emph{유도 대수기하학 I--V}
\cite{dag1}, \cite{dag2}, \cite{dag3}, \cite{dag4}, \cite{dag5}
\item To\"en: \emph{고차 및 유도 스택: 전역적 개관}
\cite{toen_higher}
\item To\"en과 Vezzosi: \emph{호모토피 대수기하학 I, II}
\cite{hag1}, \cite{hag2}
\end{itemize}
